%%%%%%%%%%%%%%%%%%%%%%%%%%%%%%%%%%%%%%%%%%%%%%%%%%%%%%%%%%%%%%%%%%%%%%%%%%%%%%%%
%% GECCO / SIGEVO submission — ACM acmart class, sigconf format.
%% Vorgaben: https://gecco-2025.sigevo.org/Paper-Submission-Instructions
%% Kompilieren: pdflatex main; bibtex main; pdflatex main; pdflatex main
%%%%%%%%%%%%%%%%%%%%%%%%%%%%%%%%%%%%%%%%%%%%%%%%%%%%%%%%%%%%%%%%%%%%%%%%%%%%%%%%
\documentclass[sigconf]{acmart}

%% ---------------------------------------------------------------------------
%% Für den anonymen Review-Prozess (Double-Blind) diese Zeile aktivieren.
%% Für die finale Version (Camera-Ready) auskommentieren.
%% \settopmatter{printacmref=false}   % ACM-Referenzblock ausblenden (Kurs)
%% ---------------------------------------------------------------------------

%% Sprache: Deutsch. babel für Silbentrennung/Anführungszeichen.
\usepackage[ngerman]{babel}

%% babel benennt die Abstract-Überschrift zu "Zusammenfassung" um -> zurück auf "Abstract".
\addto\captionsngerman{\renewcommand{\abstractname}{Abstract}}

%% Kein ACM-Copyright/DOI-Block für eine Uni-Ausarbeitung.
\setcopyright{none}
\settopmatter{printacmref=false}   % ACM-Reference-Format-Block entfernen
\renewcommand\footnotetextcopyrightpermission[1]{} % Fußzeilen-Copyright weg
\pagestyle{plain}                  % Seitenzahlen, keine ACM-Kopfzeile

\begin{document}

\title{Wie viel Wahrnehmung braucht ein NEAT-Agent?}
\subtitle{Einfluss von Zielrichtung und Sichtradius auf das Lösen von Labyrinthen}

%% ---------------------------------------------------------------------------
%% Autoren
%% ---------------------------------------------------------------------------
\author{Tobias Bandt}
\affiliation{%
  \institution{HAW Hamburg}
  \city{Hamburg}
  \country{Deutschland}
}
\email{mail@tobiasbandt.de}

%% Kopfzeilen-Kurztitel für Autoren
\renewcommand{\shortauthors}{Bandt}

%% ---------------------------------------------------------------------------
%% Abstract
%% ---------------------------------------------------------------------------
\begin{abstract}
  Abstract-Text here
\end{abstract}

%% CCS-Concepts optional. Für den Kurs weglassen oder minimal halten.
% \begin{CCSXML} ... \end{CCSXML}
% \ccsdesc[500]{Computing methodologies~Genetic algorithms}

\keywords{Neuroevolution, NEAT, Labyrinth, Wahrnehmungsradius, Evolutionäre Algorithmen}

\maketitle

%% ===========================================================================
\section{Einführung}
Kontext, Kurs erklären, praktikum.

\section{Ausgangszustand}
Aktueller Code: NEAT-Setup, 8-Zellen-Input, Fitness, Konfiguration.

\section{Vorgehen}
Festlegen was getestet werden soll und erwarteter Effekt notieren

\section{Durchführung}
Code-Diffs, code erklären, Metriken, Seeds/Mazes.

\section{Ergebnisse}
Nur Fakten/Messwerte, Diagramme pro Konfiguration.

\section{Diskussion}
Interpretation, Hypothesen bestätigt/verworfen, eigene Einschätzung.

\section{Fazit}
Rückblick auf das Projekt.

%% ===========================================================================
%% Literatur: ACM-Reference-Format (bibtex). Datei references.bib.
%\bibliographystyle{ACM-Reference-Format}
%\bibliography{references}

\end{document}
